% ===== Section 1: A rose by any other name ... (smell and taste, ~16 min) =====
\topicopener{nose}{Հոտ ու համ}{Վարդը, ինչ անուն էլ տաս...}{Ինչու է հոտառությունն ավելի ուժեղ, քան համը}{hungry}

\begin{frame}{Քանի՞ զգայարան ունենք}
  \begin{columns}[T]
    \begin{column}{0.46\textwidth}
      \begin{chatbox}
        Քանի՞ զգայարան ունի մարդը։ Գրիր չաթում \textbf{մեկ թիվ}։
      \end{chatbox}
      \vspace{1mm}
      \onslide<3->{%
      {\blackfont Ու կան «լրացուցիչները»}\par\smallskip
      \chip{hand}{պրոպրիոցեպցիա}\par\smallskip
      \chip{ear}{հավասարակշռություն}\par\smallskip
      \chip{doughnut}{ախորժակ}\ {\footnotesize\color{plum!70!cream}(վեբինար 1)}}
    \end{column}
    \begin{column}{0.5\textwidth}
      \onslide<2->{%
      {\blackfont Հինգ «իսկական» զգայարանները}\par\smallskip
      \begin{tikzpicture}[row/.style={draw=plum, line width=0.7pt, fill=cardwhite, rounded corners=2mm,
          text width=6.1cm, inner sep=3pt, anchor=west, font=\small, text=plum}]
        \foreach \i/\e/\a/\b in {0/eye/տեսողություն/ophthalmoception, 1/ear/լսողություն/audioception,
            2/tongue/համ/gustaoception, 3/nose/հոտ/olfacoception, 4/hand/հպում/tactioception}{
          \node[row] at (0,-\i*0.56) {\emoji[1.05em]{\e}\ \textbf{\a}\hfill{\scriptsize\color{plum!65!cream}\b}};}
      \end{tikzpicture}\par\smallskip
      {\footnotesize Գիտական անունները գիտնականներն էլ գրեթե չեն օգտագործում. ժամանակ են խնայում \emoji{hourglass}}}
    \end{column}
  \end{columns}
\end{frame}

\begin{frame}{Հոտառությունը՝ ամենահին զգայարանը}
  \begin{columns}[T]
    \begin{column}{0.5\textwidth}
      \begin{itemize}
        \item Հոտառությունը, ենթադրաբար, էվոլյուցիայի ընթացքում \textbf{առաջինն է} առաջացել
        \item Քիթը ուղիղ աչքերի տակ է, բայց հոտառությունը հաճախ աչքաթող ենք անում
        \item<2-> Մոր շնչած մասնիկները հասնում են պտղաջրին, ու պտուղը \textbf{զգում է մոր զգացած հոտերը}
      \end{itemize}
    \end{column}
    \begin{column}{0.46\textwidth}
      \onslide<3->{%
      \begin{scienceupdate}
        Գրքում՝ արգանդում \textbf{առաջինը} հոտառությունն է զարգանում։ Սովորաբար առաջինը համարում են \textbf{հպումը}։ Բայց որ պտուղը սովորում է մոր կերածի հոտը, հաստատված է։
      \end{scienceupdate}}
    \end{column}
  \end{columns}
  \onslide<3->{\source{Lecanuet \& Schaal, 1996; Schaal et al., Chem. Senses, 2000}}
\end{frame}

\begin{frame}{10 հազար, թե՞ 1 տրիլիոն հոտ}
  \begin{columns}[T]
    \begin{column}{0.52\textwidth}
      \centering
      % How many odours can humans tell apart? Log scale (one tick = x1000): the 1920s estimate of 10,000
% vs the ~1 trillion of Bushdid et al. 2014. The second bar appears on slide 2 (use inside a frame).
\begin{tikzpicture}[x=0.56cm]
  \draw[plum, line width=0.8pt] (0,-0.15) -- (12.4,-0.15);
  \foreach \p/\t in {0/1, 3/հազար, 6/միլիոն, 9/միլիարդ, 12/տրիլիոն}{
    \draw[plum, line width=0.8pt] (\p,-0.15) -- (\p,-0.27);
    \node[font=\scriptsize, text=plum, below, inner sep=2pt] at (\p,-0.25) {\t};}
  \fill[amBlue] (0,1.3) rectangle (4,1.85);
  \node[anchor=west, font=\blackfont\small, text=plum] at (4.1,1.575) {10\,000};
  \node[anchor=south west, font=\scriptsize, text=plum, inner sep=1pt] at (0,1.87) {1920-ականներ, տեսական գնահատական};
  \fill[amOrange, visible on=<2->] (0,0.1) rectangle (12,0.65);
  \node[visible on=<2->, anchor=east, font=\blackfont\small, text=plum] at (11.9,0.375) {1\,000\,000\,000\,000};
  \node[visible on=<2->, anchor=south west, font=\scriptsize, text=plum, inner sep=1pt] at (0,0.67) {2014, Bushdid et al., Science};
  \node[font=\scriptsize, text=plum!70!cream] at (6,-0.85) {ամեն գծիկ՝ 1000 անգամ ավելի};
\end{tikzpicture}
\par\medskip
      \onslide<4->{%
      \begin{scienceupdate}
        Բըրնեթն ինքն էլ է զգուշացնում. իրական թիվը \textbf{դեռ հայտնի չէ}։
      \end{scienceupdate}}
    \end{column}
    \begin{column}{0.44\textwidth}
      \begin{itemize}
        \item 1920-ականներ. \textbf{10\,000}՝ առանց ստուգելու
        \item<2-> 2014. մարդիկ տարբերեցին իրար շատ նման հոտերի խառնուրդներ. \textbf{1 տրիլիոն}
      \end{itemize}
      \onslide<3->{%
      \begin{bookquote}
        Կարծես փոշեկուլի պահարանում խլուրդ-մարդկանց քաղաք գտնես։
      \end{bookquote}}
    \end{column}
  \end{columns}
  \onslide<4->{\source{Gerkin \& Castro, eLife, 2015}}
\end{frame}

\begin{frame}{Ինչպես ենք հոտ զգում}
  \centering
  % The smell pathway as in the book: molecules -> olfactory neurons in the upper nose -> olfactory bulb
% -> olfactory nucleus and piriform cortex. Reveals in 4 steps (use inside a frame).
\begin{tikzpicture}[st/.style={draw=plum, line width=0.8pt, rounded corners=2.5mm, align=center, font=\small,
    text=plum, inner sep=4pt, text width=2.75cm, minimum height=2.3cm}]
  \node[st, fill=brainpink!35!cardwhite] (a) at (0,0) {\bigemoji[0.75cm]{rose}\\[1mm] Մոլեկուլներ օդում};
  \node[st, fill=peach!45!cardwhite, visible on=<2->] (b) at (3.5,0)
    {\bigemoji[0.75cm]{nose}\\[1mm] \textbf{Հոտառական նեյրոններ}\\ քթի վերին մասում};
  \node[st, fill=lavender!45!cardwhite, visible on=<3->] (c) at (7.0,0)
    {\bigemoji[0.75cm]{brain}\\[1mm] \textbf{Հոտառական սոխուկ}\\ հավաքում ու դասավորում է};
  \node[st, fill=mint!45!cardwhite, visible on=<4->] (d) at (10.5,0)
    {\bigemoji[0.75cm]{sparkles}\\[1mm] \textbf{Կեղև}\\ (piriform cortex)\\ «Հոտ եմ զգում»};
  \draw[warr, visible on=<2->] (a) -- (b);
  \draw[warr, visible on=<3->] (b) -- (c);
  \draw[warr, visible on=<4->] (c) -- (d);
\end{tikzpicture}
\par\medskip
  \onslide<5->{%
  \begin{card}[width=0.86\textwidth, colback=butter!40!cream]
    \centering Հոտառական նյարդը \textbf{12 գանգուղեղային նյարդերից} (cranial nerves) առաջինն է։ Երկրորդը տեսողականն է։
  \end{card}}
\end{frame}

\begin{frame}{Քթի երկու գաղտնիք}
  \begin{columns}[T]
    \begin{column}{0.47\textwidth}
      \begin{card}[colback=mint!35!cardwhite]
        \emoji[1.4em]{sparkles}\ {\blackfont Նյարդային համակարգի Վուլվերինը}\par\smallskip
        Հոտառական նեյրոնները շփվում են դրսի աշխարհի հետ ու մաշվում են։ Դրա համար նրանք մարդու այն քիչ նեյրոններից են, որոնք \textbf{վերականգնվում են}։
      \end{card}
    \end{column}
    \begin{column}{0.47\textwidth}
      \onslide<2->{%
      \begin{card}[colback=lavender!40!cardwhite]
        \emoji[1.4em]{dna}\ {\blackfont Հազար գեն}\par\smallskip
        Աքսելն ու Բաքը (1991) գտան հոտառական ընկալիչների մոտ \textbf{1000 գեն}՝ մեր գեների մոտ \textbf{3\%}-ը։ 2004-ին՝ Նոբելյան մրցանակ։
      \end{card}}
    \end{column}
  \end{columns}
  \vspace{2mm}
  \onslide<3->{{\small Վերականգնվող նեյրոնները ուսումնասիրում են, որ նույնը «սովորեցնեն» ողնուղեղի վնասված նեյրոններին։}\par}
  \onslide<2->{\source{Buck \& Axel, Cell, 1991; Նոբելյան կոմիտե, 2004}}
\end{frame}

\begin{frame}{Հոտն ու հիշողությունը հարևաններ են}
  \begin{columns}[c]
    \begin{column}{0.5\textwidth}
      \centering
      \begin{tikzpicture}
        \node[wbox=brainpink!40!cardwhite, text width=2.4cm] (ob) at (0,0) {\emoji{nose}\\ \textbf{Հոտառական սոխուկ}};
        \node[wbox=lavender!45!cardwhite, text width=2.4cm] (hc) at (3.5,1.05) {\emoji{books}\ \textbf{Հիպոկամպ}\\ հիշողություն};
        \node[wbox=peach!50!cardwhite, text width=2.4cm] (am) at (3.5,-1.05) {\emoji{heart}\ \textbf{Նշաձև մարմին}\\ զգացմունքներ};
        \draw[{Stealth[length=2.4mm]}-{Stealth[length=2.4mm]}, plum, line width=1.2pt] (ob) -- (hc);
        \draw[{Stealth[length=2.4mm]}-{Stealth[length=2.4mm]}, plum, line width=1.2pt] (ob) -- (am);
        \begin{scope}[on background layer]
          \draw[plum!50, dashed, line width=0.8pt, rounded corners=4mm] (-1.5,-1.95) rectangle (5.0,1.95);
        \end{scope}
        \node[font=\scriptsize\mediumfont, text=plum!80!cream, fill=cream, inner sep=2pt] at (1.75,1.95) {լիմբիկ համակարգ (limbic system)};
      \end{tikzpicture}
    \end{column}
    \begin{column}{0.46\textwidth}
      \begin{itemize}
        \item Հոտառական համակարգը հիշողության համակարգի կողքին է
        \item<2-> Բայց նրանք ոչ թե պարզապես հարևաններ են, ինչպես վեգանն ու մսագործը, այլ \textbf{կապված են}
        \item<3-> Դրա համար հոտը արթնացնում է \textbf{ամենավառ} հիշողություններն ու զգացմունքները
      \end{itemize}
    \end{column}
  \end{columns}
\end{frame}

\begin{frame}{Մի հոտ, ու դու մանկության մեջ ես}
  \begin{columns}[T]
    \begin{column}{0.46\textwidth}
      \begin{chatbox}
        Ո՞ր հոտն է քեզ միանգամից տանում մանկություն։ Գրիր չաթում։
      \end{chatbox}
      \vspace{1mm}
      {\small Բըրնեթի օրինակները՝ կիրակնօրյա ճաշը տատիկի տանը, պապիկի ծխամորճը։}
    \end{column}
    \begin{column}{0.5\textwidth}
      \onslide<2->{%
      \begin{card}
        \emoji[1.2em]{bread}\ Հաց \textbf{տեսնելը} սովորական է, թարմ հացի \textbf{հոտը}՝ հաճելի։\par\smallskip
        \emoji[1.2em]{nauseated}\ Փտած միսը տհաճ է տեսնել, բայց \textbf{հոտն} է սիրտ խառնում։
      \end{card}}
      \onslide<3->{%
      \begin{card}[colback=butter!45!cream]
        \emoji[1.2em]{moneybag}\ Խանութները հոտով տրամադրություն են ստեղծում, որ շատ ծախսես։ Բայց վանիլային պաղպաղակից թունավորվածին վանիլը չի հանգստացնի։
      \end{card}}
    \end{column}
  \end{columns}
\end{frame}

\begin{frame}{Հոտն էլ է խաբվում}
  \begin{columns}[c]
    \begin{column}{0.46\textwidth}
      \begin{itemize}
        \item Երկար ժամանակ կարծում էին, որ հոտառությունը \textbf{հնարավոր չէ խաբել}
        \item<2-> 2001. Հերցն ու ֆոն Քլեֆը (Herz \& von Clef) մարդկանց տվեցին նույն հոտը տարբեր պիտակներով
        \item<2-> «Տոնածառ» պիտակով այն հաճելի էր, «զուգարանի մաքրող միջոց» պիտակով՝ տհաճ
        \item<3-> Բըրնեթ՝ «սա կատակ չէ, իրական օրինակ է»
      \end{itemize}
    \end{column}
    \begin{column}{0.5\textwidth}
      \centering
      \onslide<4->{%
      \begin{meme}[6.6cm]{spiderman_pointing.jpg}
        \memecap{0.3}{0.88}{«պարմեզան»}
        \memecap{0.72}{0.88}{«փսխուք»}
        \memecap[font=\blackfont\footnotesize]{0.5}{0.1}{նույն շիշը}
      \end{meme}}
    \end{column}
  \end{columns}
  \onslide<4->{\source{Herz \& von Clef, Perception, 2001}}
\end{frame}

\begin{frame}{Իլյուզիա, թե հալյուցինացիա}
  \begin{columns}[T]
    \begin{column}{0.47\textwidth}
      \begin{card}[colback=mint!35!cardwhite]
        \emoji[1.4em]{magicwand}\ {\blackfont Իլյուզիա} (illusion)\par\smallskip
        Ինչ-որ բան \textbf{կա}, բայց ուղեղը այն սխալ է մեկնաբանում։ Սովորական «քմահաճույք» է։\par\smallskip
        Հոտը ուղեղին քիչ բան է ասում՝ «կա» կամ «չկա»։ Եթե ուղեղը ազդանշանները խառնի, ինչպե՞ս կիմանաս։
      \end{card}
    \end{column}
    \begin{column}{0.47\textwidth}
      \onslide<2->{%
      \begin{card}[colback=brainpink!30!cardwhite]
        \emoji[1.4em]{fire}\ {\blackfont Հալյուցինացիա}\par\smallskip
        Զգում ես մի բան, որը \textbf{չկա}։ Իրական անսարքություն է։\par\smallskip
        Հաճախ «այրվածի» ֆանտոմ հոտ է՝ տոստ, ռետին, մազ։ Կարող է կապված լինել էպիլեպսիայի, ուռուցքի, ինսուլտի հետ։
      \end{card}}
    \end{column}
  \end{columns}
  \vspace{2mm}
  \onslide<3->{\centering\small Այս տարբերությունը պետք կգա ամբողջ գլխի ընթացքում։\par}
\end{frame}

\begin{frame}{Համը՝ երկրորդ քիմիական զգայարանը}
  \begin{columns}[T]
    \begin{column}{0.56\textwidth}
      \begin{itemize}
        \item Արձագանքում է \textbf{թքի մեջ լուծված} քիմիական նյութերին
        \item Ընկալիչները լեզվի \textbf{համային պտկիկներում} են (taste buds)
      \end{itemize}
      \vspace{1mm}
      \begin{tikzpicture}
        \foreach \x/\e/\t/\c in {0/salt/աղի/cardwhite, 1.4/candy/քաղցր/brainpink!35!cardwhite, 2.8/coffee/դառը/peach!45!cardwhite,
            4.2/lemon/թթու/butter!60!cardwhite, 5.6/cheese/ումամի/mint!45!cardwhite}{
          \node[wtile=\c] at (\x,0) {\emoji[1.5em]{\e}\\[1pt] \textbf{\t}};}
      \end{tikzpicture}\par\smallskip
      {\small Ումամին գլուտամատի՝ «մսի» համն է։ Կան նաև տտիպը (լոռամիրգ), կծուն (կոճապղպեղ), մետաղականը (...մետաղ)}
    \end{column}
    \begin{column}{0.4\textwidth}
      \onslide<2->{%
      \begin{scienceupdate}
        Գրքում՝ «պտկիկների հինգ տեսակ կա»։ Իրականում հինգ տեսակի են \textbf{ընկալիչ բջիջները}, ու ամեն պտկիկում կան բոլորը։ Լեզվի բոլոր մասերն էլ զգում են \textbf{բոլոր} համերը։
      \end{scienceupdate}}
    \end{column}
  \end{columns}
  \onslide<2->{\source{Chandrashekar et al., Nature, 2006}}
\end{frame}

\begin{frame}{Համը՝ ամենաթույլ զգայարանը}
  \begin{columns}[c]
    \begin{column}{0.62\textwidth}
      \begin{bookquote}
        Գինու գիտակը մի կում է անում ու հայտարարում՝ սա 54 տարեկան ֆրանսիական «Շիրազ» է՝ կաղնու, մշկընկույզի, նարնջի ու խոզի մսի նրբերանգներով, իսկ խաղողը ճզմել է 28-ամյա Ժակը, որի ձախ կրունկին գորտնուկ կա։
      \end{bookquote}
    \end{column}
    \begin{column}{0.34\textwidth}
      \centering
      \bigemoji[1.3cm]{wine}\par\medskip
      \onslide<2->{%
      \begin{card}[colback=peach!40!cardwhite]
        Նման ճշգրիտ քիմքը ավելի շատ \textbf{մտքի} գործ է, քան լեզվի։
      \end{card}}
    \end{column}
  \end{columns}
\end{frame}

\begin{frame}{Առասպել, թե՞ փաստ}
  \mythorfact{Գինու փորձագետը միշտ կտարբերի սպիտակ գինին կարմիրից}{ԱՌԱՍՊԵԼ}{pinkdark}{%
    \begin{card}[width=0.8\textwidth, colback=brainpink!25!cardwhite]
      \emoji[1.2em]{testtube}\ Գիտակներին տվել են \textbf{կարմիր ներկած սպիտակ գինի}, ու նրանք այն նկարագրել են որպես կարմիր։
    \end{card}}
  \onslide<2->{\source{Morrot, Brochet \& Dubourdieu, Brain Lang., 2001}}
\end{frame}

\begin{frame}{Գիտակներն էլ են մարդ}
  \begin{columns}[T]
    \begin{column}{0.5\textwidth}
      \begin{itemize}
        \item Մասնագետների կարծիքները \textbf{չեն համընկնում}. մեկի համար գինին լավագույնն է, մյուսի համար՝ ճահճաջուր
        \item<2-> Հաճախ չեն տարբերում հայտնի հին գինին էժան զանգվածայինից
        \item<3-> Նրանք \textbf{չեն ստում}. իսկապես զգում են այդ համերը, բայց դրանք հիմնականում ակնկալիքից ու փորձից են
      \end{itemize}
    \end{column}
    \begin{column}{0.46\textwidth}
      \onslide<4->{%
      \begin{bookquote}
        Գիտնականները գինու գիտակների դեմ տարօրինակ ոխ չունեն։ Պարզապես քիչ մասնագիտություններ են այդքան հույս դնում համի վրա։
      \end{bookquote}}
    \end{column}
  \end{columns}
  \source{Hodgson, J. Wine Econ., 2009}
\end{frame}

\begin{frame}{Համի մեծ մասը հոտն է}
  \begin{columns}[c]
    \begin{column}{0.52\textwidth}
      \begin{itemize}
        \item Մրսած մարդիկ դժգոհում են, որ ուտելիքի համը չեն զգում
        \item<2-> Փակ աչքերով ու քթով մարդիկ չեն տարբերում \textbf{խնձորը, կարտոֆիլն ու սոխը} {\small(հիշո՞ւմ ես վեբինար 1-ի փորձը)}
        \item<3-> Օվրեյն ու Սփենսը (Auvray \& Spence). ուտելիս ուղեղը հոտը ընկալում է որպես \textbf{համ}, որովհետև զգացողությունների մեծ մասը բերանում է
      \end{itemize}
    \end{column}
    \begin{column}{0.44\textwidth}
      \centering
      % Why smell passes for taste: while you chew, odour molecules rise from the back of the mouth to the
% smell receptors at the top of the nose; the brain files them under "taste" (Auvray & Spence).
% Head in profile, facing left. Reveals on slide 3 (use inside a frame).
\begin{tikzpicture}[scale=1.15]
  \path[use as bounding box] (-1.1,-1.75) rectangle (4.25,2.35);
  % head
  \filldraw[fill=peach!35!cardwhite, draw=plum, line width=1pt]
    plot[smooth, tension=0.6] coordinates {(1.15,-1.7) (1.0,-1.05) (0.45,-0.85) (0.0,-0.72) (-0.08,-0.5)
      (0.0,-0.3) (-0.08,-0.15) (0.0,0.0) (-0.1,0.12) (0.0,0.3)}
    -- plot[smooth, tension=0.6] coordinates {(0.0,0.3) (-0.5,0.5) (-0.12,1.2) (-0.12,1.6) (0.3,2.1) (1.2,2.3)
      (2.1,2.0) (2.55,1.3) (2.55,0.45) (2.2,-0.3) (2.05,-1.0) (2.05,-1.7)} -- cycle;
  \fill[plum] (0.25,1.28) circle (0.06);
  % nasal cavity and mouth
  \draw[plum!70, line width=0.8pt, fill=cardwhite]
    plot[smooth, tension=0.7] coordinates {(-0.25,0.42) (0.25,0.75) (0.8,0.95) (1.3,0.8) (1.45,0.4) (1.45,-0.1)
      (1.25,-0.25) (0.8,-0.2) (0.2,-0.12) (0.05,-0.08)}
    -- (0.08,0.12)
    -- plot[smooth, tension=0.7] coordinates {(0.08,0.12) (0.5,0.12) (1.05,0.25) (1.15,0.55) (0.8,0.7) (0.3,0.55) (-0.18,0.32)} -- cycle;
  % smell receptors at the top of the nose
  \foreach \x in {0.55,0.72,0.89,1.06}{\fill[pinkdark] (\x,0.92) circle (0.045);}
  \node at (0.6,-0.1) {\bigemoji[0.32cm]{candy}};
  % brain says "taste"
  \node[say, callout absolute pointer={(1.3,2.15)}, text width=1.8cm, font=\handfont\small] at (2.9,2.0) {Ինչ համեղ է};
  % the route of the smell from the mouth
  \draw[-{Stealth[length=2.2mm]}, amOrange!85!black, line width=1.4pt, visible on=<3->]
    (0.9,-0.1) .. controls (1.4,-0.1) and (1.4,0.6) .. (1.0,0.82);
  \node[visible on=<3->, font=\scriptsize, text=plum, align=left, anchor=west] at (2.65,0.35)
    {հոտը քիթ է\\ բարձրանում\\ \textbf{բերանի}\\ \textbf{ներսից}};
  \draw[visible on=<3->, plum!60, line width=0.6pt] (2.65,0.35) -- (1.5,0.35);
  \node[font=\scriptsize, text=plum, anchor=east] at (-0.15,0.95) {քիթ};
  \node[font=\scriptsize, text=plum, anchor=east] at (-0.15,-0.1) {բերան};
\end{tikzpicture}

    \end{column}
  \end{columns}
\end{frame}

\begin{frame}{Փորձիր հիմա. կոնֆետ առանց քթի}
  \begin{columns}[c]
    \begin{column}{0.6\textwidth}
      \begin{tryit}
        \begin{enumerate}
          \item Մատներով ամուր \textbf{սեղմիր քիթդ}
          \item Բերանդ դիր կոնֆետ, շոկոլադի կամ մրգի կտոր ու ծամիր 10 վայրկյան
          \item Ի՞նչ ես զգում. քաղցր, թե՞ կոնկրետ համ
          \item Հիմա \textbf{բաց թող քիթդ}
        \end{enumerate}
      \end{tryit}
    \end{column}
    \begin{column}{0.36\textwidth}
      \centering
      \bigemoji[1.3cm]{candy}\par\medskip
      \onslide<2->{%
      \begin{card}[colback=mint!35!cardwhite]
        Քաղցրը լեզուն է զգում։ Իսկ ելակը, կիտրոնը կամ անանուխը «գալիս են» միայն քիթը բացելուց հետո։
      \end{card}}
    \end{column}
  \end{columns}
\end{frame}

\begin{frame}{Խորհուրդ վատ խոհարարներին}
  \begin{columns}[c]
    \begin{column}{0.58\textwidth}
      \begin{bookquote}
        Եթե վատ խոհարար ես, միևնույն է, կարող ես ընթրիքի հյուրեր հրավիրել, եթե նրանք սարսափելի մրսած են ու պատրաստ են նստել մթության մեջ։
      \end{bookquote}
      \vspace{2mm}
      \onslide<2->{%
      \begin{card}[colback=butter!45!cream]
        \centering Հոտը \textbf{թերագնահատված} է, համը՝ \textbf{գերագնահատված}։ Ու ուղեղը անընդհատ խառնում է դրանք։
      \end{card}}
    \end{column}
    \begin{column}{0.36\textwidth}
      \centering
      \mascot[3cm]{hungry}\par
      \bigemoji[0.8cm]{forkknife}\ \bigemoji[0.8cm]{moon}
    \end{column}
  \end{columns}
\end{frame}
